\documentclass[11pt]{article}
\usepackage[T1]{fontenc}
\usepackage{lmodern,amsmath,amssymb,amsthm,mathtools}
\usepackage[margin=1in]{geometry}
\usepackage{microtype,enumitem,booktabs,longtable,array,calc}
\usepackage{fancyvrb,fvextra}
\usepackage{graphicx}
\usepackage[hidelinks,unicode]{hyperref}
\usepackage{xurl}
\hypersetup{pdftitle={Tuza's Inequality and Equality Characterization for 5-Degenerate Graphs},pdfauthor={AI Math Experiments},pdfsubject={Triangle edge packing, covering, equality, and formal verification},pdfkeywords={Tuza, five-degenerate graphs, triangle packing, Lean},pdfcreator={LaTeX with Tectonic},pdfproducer={Tectonic}}
\setlength{\emergencystretch}{3em}
\setlength{\parskip}{3pt}
\setcounter{secnumdepth}{-1}
\setcounter{tocdepth}{1}
\setlist{itemsep=2pt,topsep=4pt}
\providecommand{\tightlist}{\setlength{\itemsep}{0pt}\setlength{\parskip}{0pt}}
\DefineVerbatimEnvironment{Highlighting}{Verbatim}{fontsize=\small,breaklines=true,commandchars=\\\{\}}
\newenvironment{Shaded}{}{}
\newcommand{\KeywordTok}[1]{#1}\newcommand{\NormalTok}[1]{#1}
\newcommand{\OtherTok}[1]{#1}\newcommand{\OperatorTok}[1]{#1}
\newcommand{\FunctionTok}[1]{#1}\newcommand{\DataTypeTok}[1]{#1}
\newcommand{\BuiltInTok}[1]{#1}\newcommand{\ControlFlowTok}[1]{#1}
\newcommand{\StringTok}[1]{#1}\newcommand{\DecValTok}[1]{#1}
\newcommand{\CommentTok}[1]{#1}\newcommand{\PreprocessorTok}[1]{#1}
\title{Tuza's Inequality and Equality Characterization\\for 5-Degenerate Graphs}
\author{AI Math Experiments}
\date{30 October 2026}
\begin{document}
\maketitle
\begin{abstract}
For a finite simple graph \(G\), let \(\nu(G)\) be the maximum number of
pairwise edge-disjoint triangles and let \(\tau(G)\) be the minimum
number of graph edges meeting every triangle. We prove
\(\tau(G)\le 2\nu(G)\) for every 5-degenerate graph. Equality holds
precisely when every nonempty triangle component, where adjacency means
sharing a graph edge, has edge support \(K_4\) or \(K_5\). Triangle-free
background edges and clique blocks sharing vertices are allowed; graph
order and maximum degree are unrestricted. The ordinary proof extends
local reducibility arguments to minimal failures of the combined
inequality and equality statement. A second degeneracy argument forces
an eight-vertex core; a symbolic compatibility argument handles
arbitrary outside attachments and saves one edge in a final cover. The
exact combined theorem is checked in Lean 4.33.1 with pinned Mathlib,
with all 280 mathematical modules rebuilt and replayed by the standard
Lean kernel. The equality classification is the principal contribution
proposed here. Historical priority and external specialist peer review
are not claimed.

\end{abstract}
\section{Introduction and precise
contribution}\label{introduction-and-precise-contribution}

Tuza's conjecture asks whether \(\tau(G)\le2\nu(G)\) holds for every
finite simple graph. This paper treats the subgraph-closed class of
5-degenerate graphs and proves an exact equality characterization within
that class.

\par\medskip\noindent\textbf{Theorem 1 (Main theorem).} Every finite simple 5-degenerate
graph \(G\) satisfies

\[\tau(G)\le2\nu(G).\]

Equality holds if and only if every nonempty shared-edge triangle
component has edge support \(K_4\) or \(K_5\).

A graph is \textbf{5-degenerate} if every nonempty vertex-induced
subgraph has a vertex of degree at most five in that subgraph.
Equivalently, every nonempty subgraph has such a vertex: deleting edges
cannot increase degrees. This is a degeneracy condition, with no bound
on original maximum degree. It does not imply maximum average degree
less than seven. The theorem includes the empty graph and every
triangle-free graph: both parameters vanish and the component condition
is vacuous.

Two triangles are adjacent if they share an \textbf{edge}. Their
connected components under chains of this adjacency are the
\emph{triangle components}. The edge support of a component is the union
of its triangle edges. This differs from ordinary graph connectedness:
distinct triangle components may share vertices, and edges lying in no
triangle affect neither parameter. All triangle families used below are
the full families of their graphs unless a subfamily is expressly named.
A string \(abc\) denotes the triangle on distinct vertices \(a,b,c\),
and \(ab\) its graph edge. We write \(N(v)\) and \(N[v]\) for open and
closed neighborhoods, \(d(v)\) for degree, and \(G[N(v)]\) for a vertex
link.

\textbf{Equivalent block formulation.} The component condition is
equivalent to the existence of a finite family of pairwise edge-disjoint
actual \(K_4\) and \(K_5\) subgraphs containing every graph triangle. To
prove this, component supports give such a family directly. Conversely,
assign each triangle to a containing block. Triangles in different
blocks cannot share an edge, by edge disjointness. Within \(K_4\) or
\(K_5\) its triangles are connected by shared edges: two triangles
sharing only a vertex are joined by a triangle using that vertex and one
other vertex from each. Every block edge belongs to a triangle. Thus a
triangle component consists exactly of one block's full triangle family,
and its literal edge support is that clique. This argument checks the
literal component condition rather than assuming it from a supplied
decomposition. Since packing and covering numbers add over edge-disjoint
triangle supports, the elementary parameters in Lemma 2 prove the
equality converse for arbitrary graphs.

The proof's starting framework is due to Puleo {[}1, Definition 2.1 and
Lemma 2.2{]}. His robustness and local low-degree results {[}1, Lemmas
2.6--2.7 and Propositions 5.1--5.2{]} are predecessors of Sections 3--6;
the four/eight reducing pair in Section 6.2 is the witness in his
Proposition 5.2, Case 1. Botler, Fernandes and Gutiérrez {[}2, Lemma
3.3{]} provide the direct predecessor of Proposition 8 for irreducible
robust graphs. The extra requirement here is that balanced reductions
need not contradict a minimal failure of the \textbf{combined}
statement. We therefore derive exact residual equality, classify its
clique supports, and prove compatibility and strict improvements at each
balanced case. These extensions, especially arbitrary outside-page
compatibility and the thirteen-versus-fourteen cover saving, remain
explicit below. Appendix A gives the shorter credited inequality
deduction, for which balanced reductions suffice.

Cui, Haxell and Ma {[}3, Theorem 1.2 and Section 2{]} characterize
planar equality using edge-disjoint \(K_4\) copies containing every
triangle, with triangle-free background edges allowed. Theorem 1
recovers their known case because planar graphs are 5-degenerate and
contain no \(K_5\). Their theorem is not claimed as a new consequence in
historical terms. Puleo's main sparsity theorem assumes maximum average
degree less than seven; we use his local framework rather than an
unsupported implication from degeneracy. Botler--Fernandes--Gutiérrez
prove a treewidth theorem; their local lemma is the antecedent used
here, not an invented degeneracy version of their main result.

The project developed the research strategy, mathematics, code, internal
reviews and formalization through AI systems. The human supplied goals,
budgets, coordination and publication authorization, and is not named as
a mathematical author. The author designation follows the project's
existing \textbf{AI Math Experiments} convention. Internal AI review is
not external human peer review. A bounded primary-source attribution and
literature-position check was carried out; historical priority remains
unestablished. No first-ever claim or general-Tuza resolution is made.

The earlier
\href{https://github.com/pappa211/ai-math-experiments/blob/main/papers/tuza-triangle-restoration/main.md}{triangle-restoration paper}
remains a separate canonical work with its own conditional
formal-verification boundary. Its restoration, odd-cycle-transversal and
density results are not retroactively claimed to have established this
theorem. General Tuza, six-degenerate graphs, OCT-five, unrestricted
two-triangle restoration and the fixed-\(c<3\) density extension are not
established here.

\subsection{Reading the proof}\label{reading-the-proof}

Sections 1--5 establish clique compatibility, tight criticality,
robustness, and the two possible degree-five links. Sections 6--7 prove
all degree-five adjacency, neighbor and overlap restrictions for a
minimal combined obstruction. Section 8 applies degeneracy a second
time, reconstructs the final core, and excludes arbitrary outside
attachments. Section 9 closes the induction. Section 10 states the exact
formal target and verification boundary; appendices give the short
inequality proof, a formally used local proof simplification, and
primary references. Every ordinary proof needed for Theorem 1 is
contained here.

For a local vertex set \(W\), an \textbf{external page} is a triangle
with base edge in \(W\) and apex outside the specified core; its
apex-to-base edges are its spokes. Different pages may share an apex. A
\emph{surviving footprint} means only edges actually remaining in the
explicitly named residual graph. No bound is imposed on outside vertices
or edges. Partial spoke deletion at a closed link component removes just
its spokes, retaining all other neighbors; closure ensures that no
unlisted hub triangle crosses to another link component.

\section{1. Elementary block facts}\label{elementary-block-facts}

\par\medskip\noindent\textbf{Lemma 2 (Clique parameters and compatibility).} The graphs
\(K_4\) and \(K_5\) have packing/covering numbers (1,2) and (2,4). A
maximum \(K_4\) packing avoids a star; a maximum \(K_5\) packing avoids
any two prescribed edges. Every edge belongs to some minimum cover.

\emph{Proof.}

\(K_4\) has \((\nu,\tau)=(1,2)\). \(K_5\) has \((\nu,\tau)=(2,4)\): two
disjoint triangles exist; three disjoint triangles would use nine edges
and leave one edge, impossible because every degree in a union of
edge-disjoint triangles is even whereas removing one edge from \(K_5\)
creates two odd degrees. Each edge of \(K_5\) meets three of its ten
triangles, so three cover edges cannot cover them all; four edges inside
a vertex bipartition of sizes two and three do cover all triangles.

A \(K_4\) maximum packing can avoid any star of forbidden edges by
choosing the face opposite the star center. A \(K_5\) maximum packing
can avoid any two edges. Up to relabeling, for forbidden adjacent edges
12,13 choose 145,234; for forbidden disjoint edges 12,34 choose 135,245.
Each displayed pair is edge-disjoint. Zero or one forbidden edge follows
by enlarging the forbidden set.

In either block there is a minimum cover containing any specified edge.
In \(K_4\) take it and its opposite edge. In \(K_5\) take the four edges
inside a \(2+3\) vertex partition whose two-vertex class is that edge's
endpoints. Independent choices combine across triangle components, whose
edge supports are disjoint.

We also use the elementary intersecting-triangle lemma: a pairwise
edge-intersecting triangle family has a two-edge cover; if it has no
common edge, all its triangles are faces of a single \(K_4\). Indeed,
two distinct triangles can be written \(abc\),\(abd\). A triangle
avoiding \(ab\) must share another edge with each, hence is \(acd\) or
\(bcd\). Once such a triangle is present, every further triangle lies
among those four faces. The pair \(ab\),\(cd\) covers the family.
Families with a common edge, or a single triangle, are immediate.

\section{2. Minimal-obstruction framework in a hereditary
class}\label{minimal-obstruction-framework-in-a-hereditary-class}

\par\medskip\noindent\textbf{Lemma 3 (Tight criticality).} An edge-minimal failure of the
combined statement in a subgraph-closed class, after removing isolated
vertices, has one nonempty triangle component, no triangle-free edges,
and \(\tau(G)=2\nu(G)\). Every edge deletion preserves \(\nu\) and
decreases \(\tau\) by one.

\emph{Proof.}

Call \(S\) the assertion that Tuza's inequality holds and equality
implies the \(K_4\)/\(K_5\) component condition. The converse has
already been proved in the introduction. A failure therefore means
either \(\tau(G)>2\nu(G)\), or \(\tau(G)=2\nu(G)\) with an unclassified
component. Let \(C\) be any class closed under taking subgraphs, and
suppose \(G\) is a failure of \(S\) in \(C\) with as few edges as
possible. Discard isolated vertices. Then every proper edge-subgraph
satisfies \(S\).

There are no triangle-free edges: deleting one preserves the triangle
family and the failure. There is only one nonempty triangle component:
otherwise apply \(S\) to each component support, sum the inequalities,
and observe that equality of the sum forces equality and the block
classification in every component. Thus \(G\) has connected nonempty
triangle support.

Write \(h=\nu(G)\). A failure of \(S\) has \(\tau(G)\ge 2h\). If
\(\nu(G-f)=h-1\) for an edge \(f\), induction gives
\(\tau(G)\le \tau(G-f)+1\le 2h-1\), impossible. Hence every edge
deletion preserves \(\nu\).

If \(\tau(G-f)=2h\) for \(f=xy\), then \(G-f\) has only \(K_4\)/\(K_5\)
triangle components. Choose a triangle \(xyw\) of \(G\). The wedge edges
\(xw\),\(yw\) cannot lie in the same residual clique block, since \(xy\)
is absent. Choose in each block a maximum packing avoiding its possible
single wedge edge and append \(xyw\). This gives \(h+1\) triangles,
contradiction. Therefore \(\tau(G-f)\le 2h-1\) for every edge \(f\).
Cover deletion changes \(\tau\) by at most one, so

\[\tau(G)=2h; \nu(G-f)=h; \tau(G-f)=2h-1 \quad\text{for every edge }f\]

Thus every minimal \(S\)-obstruction in a hereditary class is
tight-critical. This argument proves the reduction; it does not presume
global Tuza.

Whenever a local deletion produces \(K\) and supplies \(r\) triangles
disjoint from every \(K\)-packing and a cover increment of size \(2r\),
equality forces

\[\nu(K)=h-r, \tau(K)=2(h-r),\]

and the inductive hypothesis classifies \(K\) into \(K_4\)/\(K_5\)
triangle components. We will always give the actual local packing and
cover increment.

\section{3. An edge cannot belong to just one
triangle}\label{an-edge-cannot-belong-to-just-one-triangle}

\par\medskip\noindent\textbf{Lemma 4 (Triangle codegree).} Every edge of a minimal combined
obstruction belongs to at least two triangles.

\emph{Proof.}

Suppose \(uv\) belongs only to \(uvw\). Delete \(uw\),\(vw\) to get
\(K\). The edge \(uv\) is triangle-free in \(K\), so every \(K\)-packing
can be appended by \(uvw\). A \(K\)-cover plus \(uw\),\(vw\) covers
\(G\). Hence \(\nu(K)=h-1\), \(\tau(K)=2h-2\) and \(K\) has the block
classification. Each block meets \(U_0=\{u,v,w\}\) at most once:
\(uw\),\(vw\) are absent, and \(uv\) is triangle-free in \(K\).

Let \(F\) be the full family of triangles lost in this deletion. They
all contain \(w\). If \(F\) is intersecting, the preceding lemma says
that \(F\) has a common edge or is contained in \(K_4\) faces. In the
latter case, since \(uvw\) belongs to \(F\) and all \(F\)-triangles
contain \(w\), absence of a common edge forces \(wuv\),\(wux\),\(wvx\)
for a fourth vertex \(x\). Their edges force the triangle \(uvx\),
contradicting uniqueness of \(uvw\). Thus \(F\) has a one-edge cover,
giving \(\tau(G)\le \tau(K)+1<2h\).

If \(F\) is not intersecting, take two disjoint members. A \(K_4\) block
has a face avoiding its possible \(U_0\) vertex, hence avoiding both
members. Each member uses at most one edge of a given \(K_5\) block, so
that block can avoid the at most two forbidden edges. These maximum
block packings combine with the two \(F\)-triangles to give \(h+1\)
triangles, contradiction. Therefore every \(G\)-edge has triangle
codegree at least two. Equivalently every vertex in every link
\(G[N(v)]\) has degree at least two.

\section{4. Every link component has at least five
vertices}\label{every-link-component-has-at-least-five-vertices}

\par\medskip\noindent\textbf{Lemma 5 (Robustness).} Every connected component of every vertex
link in a minimal combined obstruction has at least five vertices; in
particular the obstruction \(G\) has minimum degree at least five.

\emph{Proof.}

Let \(C_0\) be a connected component of the link at \(v\). In the
following cases, deleting \(v\) means deleting just the spokes into
\(C_0\); other neighbors of \(v\) remain. There is no triangle through a
deleted spoke with its third vertex outside \(C_0\), because that would
join link components. All listed covers therefore account for every lost
triangle. External rim pages have apex outside \(C_0\cup \{v\}\).

\subsection{\texorpdfstring{4.1 Link component
\(K_3\)}{4.1 Link component K\_3}}\label{link-component-k_3}

Write \(C_0=abc\). Delete its three \(v\)-spokes and \(ab\). The
triangle \(vab\) appends to any residual \(K\)-packing, and
\(\{ab,vc\}\) is a cover increment, so \(K\) has equality blocks and
\(\nu(K)=h-1\).

An external \(abx\) combines with \(vbc\). Their edges present in \(K\)
are \(bc\),\(ax\),\(bx\). Since \(ab\) is deleted, no block contains
both \(a\) and \(b\); the only possible two-edge overlap is
\(bc\),\(bx\), a star at \(b\). Block avoidance yields \(h+1\)
triangles, contradiction. By symmetry no rim edge of \(abc\) has an
external triangle. Thus the \(K_4\) on \(v\),\(a\),\(b\),\(c\) is an
isolated triangle component. Connectedness makes it the whole triangle
support, which satisfies \(S\), contradiction.

\subsection{\texorpdfstring{4.2 Link component
\(K_4\)}{4.2 Link component K\_4}}\label{link-component-k_4}

Write \(C_0=abcd\). Delete its \(v\)-spokes and \(ab\),\(cd\). The
packing \(vab\),\(vcd\) appends, and \(\{ab,cd,va,vb\}\) is a four-edge
cover increment. Hence \(\nu(K)=h-2\) and \(K\) has equality blocks.

An external \(abx\) combines with \(vac\),\(vbd\). Its old forbidden
edges are \(ac\),\(bd\),\(ax\),\(bx\). Because \(ab\) is absent, a block
sees at most two of these edges, necessarily adjacent at \(a\) or \(b\).
Block avoidance gives \(h+1\) triangles. Varying the perfect matching
\(ab\),\(cd\) excludes external triangles on every rim edge. The \(K_5\)
on \(v\),\(a\),\(b\),\(c\),\(d\) is therefore an isolated triangle
component, again impossible.

\subsection{4.3 Diamond link component}\label{diamond-link-component}

Write \(C_0=K_4-ab\) on \(a\),\(b\),\(c\),\(d\). Delete its \(v\)-spokes
and \(ac\),\(bd\). Use packing \(vac\),\(vbd\) and cover increment
\(\{ac,bd,vc,vd\}\); equality gives \(\nu(K)=h-2\) and blocks.

An external \(acx\) combines with \(vad\),\(vbc\). Its old forbidden
edges are \(ad\),\(bc\),\(ax\),\(cx\). Since \(ac\) is absent in \(K\),
a block sees at most two of these, with the only possible pairs
\(ad\),\(ax\) or \(bc\),\(cx\). Block avoidance contradicts maximality.
By symmetry, including use of the other matching \(ad\),\(bc\), no
external triangle uses \(ac\),\(ad\),\(bc\),\(bd\). Returning to \(K\)
above, a \(K\)-cover plus \(\{vc,vd,cd\}\) covers \(G\): \(vc\),\(vd\)
cover all \(v\)-triangles, and the only other lost triangles are
\(acd\),\(bcd\), hit by \(cd\). This gives \(2h-1\) cover edges,
contradiction.

\subsection{4.4 Four-cycle link
component}\label{four-cycle-link-component}

Let \(C_0\) be the cycle \(abcd\). Delete its \(v\)-spokes and
\(ab\),\(cd\),\(da\), retaining \(bc\). The packing \(vab\),\(vcd\)
appends; all four rim edges are a cover increment. Thus \(\nu(K)=h-2\),
\(\tau(K)=2h-4\) and \(K\) has equality blocks.

By Section 3, \(bc\) has a triangle \(bcx\) other than \(vbc\). It
survives in \(K\): \(x\) cannot be \(a\) or \(d\) since the diagonals
\(ac\),\(bd\) are absent. Hence \(bc\) lies in a residual block. Choose
a minimum \(K\)-cover containing \(bc\), and add only
\(ab\),\(cd\),\(da\). It covers \(G\) with \(2h-1\) edges,
contradiction.

Since every link has minimum degree at least two, its connected
components of order at most four are exactly
\(K_3\),\(C_4\),diamond,\(K_4\). All are excluded. Hence every link
component has order at least five, and \(\delta(G)\ge 5\).

\section{5. A 3-colorable link component cannot have five
vertices}\label{a-3-colorable-link-component-cannot-have-five-vertices}

\par\medskip\noindent\textbf{Lemma 6 (Degree-five links).} The link of every degree-five
vertex in a minimal combined obstruction is \(K_5\) or \(K_5\) minus one
edge.

\emph{Proof.}

Suppose a link component \(L\) has five vertices, minimum degree at
least two, and is 3-colorable. Its complement has maximum degree at most
two. A 3-coloring gives either a size-three color class (a triangle in
the complement) or two size-two classes (a matching of size two). There
are \textbf{eleven}, not thirteen, unlabeled graphs on five vertices of
maximum degree at most two: \(5I\), \(K_2+3I\), \(2K_2+I\), \(P_3+2I\),
\(P_3+K_2\), \(P_4+I\), \(P_5\), \(C_3+2I\), \(C_3+K_2\), \(C_4+I\), and
\(C_5\). Here \(I=K_1\) and \(+\) denotes disjoint union. The first,
second, and fourth entries have non-3-colorable complements; the other
eight are the possible complements here. Thus the possible complements
are exactly

\[2K_2+I, P_3+K_2, P_4+I, P_5, C_3+K_2, C_3+2I, C_4+I, C_5\]

For a matching \(M\) and vertex set \(Q\) with \(L-Q\) a subgraph of
\(M\), delete the local \(v\)-spokes and \(M\). The \(|M|\) triangles at
\(v\) on \(M\) append to every residual packing, while \(M\) plus the
\(v\)-spokes to \(Q\) is a cover increment. If \(|Q|<|M|\) this is
strict; if equality holds, the residual has \(K_4\)/\(K_5\) blocks and
packing \(h-|M|\).

\subsection{5.1 Four dense types}\label{four-dense-types}

In each dense case \(|M|=|Q|=2\). The listed three triangles are
mutually edge-disjoint. Their edges present in the residual are listed
as their surviving footprint. The missing and deleted edges restrict a
block's intersection as shown, so maximum block packings avoid that set
and append all three triangles, contradiction.

\textbf{Dense case 1.} Let \(L\) be \(K_{2,3}\) with sides
\(\{p,q\}\),\(\{b,c,x\}\), plus \(bc\). Take \(M=\{bc,px\}\) and
\(Q=\{p,q\}\). The three triangles are \(pbc\); \(vqb\); \(vpx\); their
surviving footprint is \(\{pb,pc,qb\}\). Possible two-edge overlap in a
block: none: \(bc\) deleted and \(pq\) absent.

\textbf{Dense case 2.} Let \(L\) be \(K_5\) minus \(ab\),\(bc\),\(cd\);
fifth vertex \(e\). Take \(M=\{bd,ce\}\) and \(Q=\{a,e\}\). The three
triangles are \(ace\); \(vad\); \(vbe\); their surviving footprint is
\(\{ac,ae,ad,be\}\). Possible two-edge overlap in a block: only
\(ae\),\(ad\).

\textbf{Dense case 3.} Let \(L\) be \(K_5\) minus
\(ab\),\(bc\),\(cd\),\(de\). Take \(M=\{ac,bd\}\) and \(Q=\{a,e\}\). The
three triangles are \(ace\); \(vad\); \(vbe\); their surviving footprint
is \(\{ae,ce,ad,be\}\). Possible two-edge overlap in a block: none.

\textbf{Dense case 4.} Let \(L\) be \(K_{2,3}\) with sides
\(\{p,q\}\),\(\{a,b,c\}\), plus \(pq\). Take \(M=\{ap,bq\}\) and
\(Q=\{p,q\}\). The three triangles are \(apq\); \(vbp\); \(vcq\); their
surviving footprint is \(\{aq,pq,bp,cq\}\). Possible two-edge overlap in
a block: only \(pq\),\(cq\).

Both allowed pairs are adjacent. For completeness, in dense case 2
\(ac\) with \(ae\) requires deleted \(ce\); \(ac\) with \(ad\) requires
absent \(cd\); \(ac\) with \(be\) requires absent \(ab\); \(ae\) with
\(be\) requires absent \(ab\); \(ad\) with \(be\) requires deleted
\(bd\). Only \(ae\),\(ad\) remains, precluding a triple. In dense case
3, the pairs respectively require deleted \(ac\), absent \(de\), absent
\(ab\), deleted \(ac\), absent \(bc\), or deleted \(bd\). In dense case
4, every pair except \(pq\),\(cq\) requires deleted \(ap\) or \(bq\), or
an absent pair among \(a\),\(b\),\(c\). Dense case 1's three pairs
require \(bc\) or \(pq\). These are all pair tests.

\subsection{5.2 Three sparse types}\label{three-sparse-types}

For a bowtie \(L=abc\cup ade\), take \(M=\{bc,de\}\), \(Q=\{a\}\). The
cover increment has size three for a compatible two-packing, immediately
strict.

For \(L=K_{2,3}\) with sides \(\{a,b\}\),\(\{c,d,e\}\), take
\(M=\{ac,bd\}\), \(Q=\{a,b\}\). An external \(acx\) combines with
\(vad\),\(vbe\). The old forbidden edges are
\(ad\),\(be\),\(ax\),\(cx\). Missing/deleted \(ac\),\(ab\),\(cd\),\(ce\)
exclude every pair except \(ad\),\(ax\), adjacent at \(a\). Hence no
external \(ac\)-page exists; symmetry handles every rim edge. There are
no triangles within \(L\), so each rim edge belongs only to its
\(v\)-triangle, contradicting Section 3.

For \(L=C_5\) with rim \(abcde\), take \(M=\{ab,cd\}\), \(Q=\{c,e\}\).
An external \(abx\) combines with \(vbc\),\(vde\). Old forbidden edges
are \(bc\),\(de\),\(ax\),\(bx\). Deleted \(ab\),\(cd\) and absent
\(ad\),\(bd\) leave only the possible pair \(bc\),\(bx\). Thus no rim
edge has an external page, by symmetry. Again each rim edge belongs to
just its \(v\)-triangle, contradiction.

\subsection{\texorpdfstring{5.3 The remaining type
\(K_{2,2,1}\)}{5.3 The remaining type K\_\{2,2,1\}}}\label{the-remaining-type-k_221}

The local four-triangle/eight-edge witness is Puleo's {[}1, Proposition
5.1, proof Case 1{]}; the residual-equality and external-page
compatibility extension is proved here.

Write the missing link edges \(ab\),\(cd\) and the fifth link vertex
\(e\). The local graph \(B\) on \(W=\{v,a,b,c,d,e\}\) is \(K_6\) minus
\(ab\),\(cd\). Delete all \(B\)-edges to \(K\). The four edge-disjoint
triangles \(vac\),\(vbd\),\(ade\),\(bce\) append to every \(K\)-packing.
All eight rim edges cover every lost triangle through \(v\) and every
other lost triangle. Hence equality forces \(\nu(K)=h-4\) and
\(\tau(K)=2h-8\), and \(K\) has equality blocks. Each block meets \(W\)
at most once, since \(W\) is independent in \(K\).

Let \(F\) be the external rim pages, triangles with a rim edge and apex
outside \(W\). The lost triangles are exactly the \(B\)-triangles and
\(F\). \(B\) has the four-edge cover \(\{va,vb,ec,ed\}\); therefore
equality forces \(\tau(F)\ge 4\).

Call the following four pairs of rim edges bad:

\[\{ac,bd\}, \{ad,bc\}, \{ae,be\}, \{ce,de\}\]

Every other pair \(f\),\(g\) is good and \(B-\{f,g\}\) has three
edge-disjoint triangles. The symmetries swapping \(a\),\(b\), swapping
\(c\),\(d\) and exchanging the two two-vertex parts reduce all good
pairs to four cases:

{\def\LTcaptype{none} % do not increment counter
{\small\begin{longtable}[]{@{}ll@{}}
\toprule\noalign{}
\(f\),\(g\) & Three triangles avoiding them \\
\midrule\noalign{}
\endhead
\bottomrule\noalign{}
\endlastfoot
\(ac\),\(ad\) & \(vae\); \(bce\); \(vbd\) \\
\(ae\),\(ce\) & \(bde\); \(vad\); \(vbc\) \\
\(ac\),\(ae\) & \(bde\); \(vad\); \(vbc\) \\
\(ac\),be & \(ade\); \(vbd\); \(vce\) \\
\end{longtable}}
}

The four cases are adjacent pairs of cycle edges, spokes at \(e\) to
different parts, an incident mixed pair, and a nonincident mixed pair.
The remaining cycle pairs and spoke pairs are precisely those listed as
bad. Thus the list covers every good pair without an appeal to
enumeration.

If \(F\) has no edge-disjoint pair on good bases, then \(\tau(F)\le 2\).
To prove this, an intersecting \(F\) already has a two-edge cover.
Otherwise choose disjoint pages on bases \(f\),\(g\), which must be bad
partners. Every page on a third base must intersect both, since a third
base is good with each bad partner. Such intersections require a common
external apex and bases meeting both \(f\) and \(g\).

For adjacent bad partners \(ae\),be (or \(ce\),\(de\)), disjointness
requires distinct apices, so no third base occurs; \(f\),\(g\) cover
\(F\). The same holds for disjoint partners with distinct apices. For
disjoint partners \(ac\),\(bd\) with common apex \(x\), a third base
meeting both is only \(ad\) or \(bc\), and its page has apex \(x\). Its
presence forces every page on \(ac\) or \(bd\) to have apex \(x\), or
else it would form a good disjoint pair with that third page. The other
possible third base likewise has apex \(x\). Thus \(F\) lies in the four
pages at \(x\) over the cycle \(ac\),\(cb\),\(bd\),\(da\), covered by
\(xa\),\(xb\). If there is no third base, \(ac\),\(bd\) cover \(F\).
This proves the claim; symmetry handles \(ad\),\(bc\).

Since \(\tau(F)\ge 4\), \(F\) has a disjoint pair on good bases
\(f\),\(g\). Combine these two pages with the displayed local
three-packing avoiding \(f\),\(g\). Each residual block meets \(W\) at
most once, and hence meets each external page in at most one old spoke.
If it meets both, the forbidden spokes share its unique \(W\) vertex. A
\(K_4\) avoids that star and a \(K_5\) avoids the at most two edges.
Maximum residual block packings thus append to all five selected
triangles, yielding \(h+1\), contradiction.

This completes exclusion of every 3-colorable five-vertex link
component.

\subsection{5.4 One additional non-3-colorable
link}\label{one-additional-non-3-colorable-link}

The link consisting of \(K_4\) on \(a\),\(b\),\(c\),\(d\) and a fifth
vertex \(x\) adjacent exactly to \(a\),\(b\) is also excluded. Take
\(M=\{cd,ax\}\) and \(Q=\{a,b\}\); \(L-Q\) is the single edge \(cd\).
Deleting the local \(v\)-spokes and \(M\) forces equality blocks with
packing \(h-2\). The triangles \(acd\),\(vax\),\(vbd\) are edge-disjoint
and have old forbidden edges \(ac\),\(ad\),\(bd\). A block cannot
contain both \(ac\),\(ad\) or both \(ac\),\(bd\) because \(cd\) is
deleted; the only possible pair is \(ad\),\(bd\), adjacent at \(d\).
Maximum block packings avoid this pair and append the three triangles,
giving \(h+1\), contradiction.

This local certificate belongs to the retained AI-developed proof
record; no external human review is asserted.

\subsection{5.5 The two remaining degree-five
links}\label{the-two-remaining-degree-five-links}

Every \(K_4\)-free graph on at most five vertices is 3-colorable.
Indeed, a vertex-minimal graph that is not 3-colorable has minimum
degree at least 3. On four vertices it is \(K_4\); on five its
complement is a matching. \(K_5\) minus zero or one edge contains
\(K_4\), while \(K_5\) minus two independent edges is \(K_{2,2,1}\) and
is 3-colorable.

Thus a degree-five link not already excluded contains \(K_4\). Its fifth
vertex has at least two neighbors there. Exactly two is the
configuration excluded in Section 5.4; three gives \(K_5\) minus one
edge, and four gives \(K_5\). Hence every degree-five vertex of a
minimal \(S\) obstruction has one of these two links.

\section{6. Degree-five vertices are independent and have no degree-six
neighbor}\label{degree-five-vertices-are-independent-and-have-no-degree-six-neighbor}

\par\medskip\noindent\textbf{Lemma 7 (Low-neighbor restriction).} Degree-five vertices form
an independent set in a minimal combined obstruction, and every neighbor
of such a vertex has degree at least seven.

\emph{Proof.}

\textbf{Attribution and extension.} The following packing/cover
accounting is Puleo's reducible-set framework (Definition 2.1, Lemma
2.2). His Proposition 5.2 is the antecedent for low-degree adjacency,
including the precise balanced witnesses in Section 6.2. For the
combined obstruction, only a strict reduction is an immediate
contradiction. Every balanced case below explicitly supplies the
additional equality argument.

We use a general reduction only with its actual witnesses. For a vertex
set \(U\), let \(P\) be \(r\) edge-disjoint triangles and \(X\) an edge
set meeting every triangle incident to \(U\) and containing every
\(P\)-edge with neither endpoint in \(U\). In \(K=G-U-X\), every
\(K\)-packing can be appended by \(P\), and every \(K\)-cover plus \(X\)
covers \(G\). Thus \(|X|<2r\) is impossible in a minimal \(S\)
obstruction. If \(|X|=2r\), equality forces \(\nu(K)=h-r\),
\(\tau(K)=2(h-r)\), and the residual clique-block classification.

The singleton matching version takes a matching \(M\) in a link and a
vertex set \(Q\) such that the link minus \(Q\) is a subgraph of \(M\).
Delete the central vertex and \(M\). The hub triangles on \(M\) give
\(|M|\) packed triangles, and \(M\) plus the spokes to \(Q\) give a
cover increment \(|M|+|Q|\). No assumption that \(Q\) and \(M\) are
disjoint is needed.

\textbf{Residual-support observation.} We will repeatedly exclude a
residual containing \(K_5\) minus one edge or minus two independent
edges. In each graph all remaining edges lie in one triangle component
with five vertices. It cannot lie in a \(K_4\) block; it cannot lie in a
\(K_5\) block because the specified missing edge was deleted from the
whole residual. This can also be seen by listing their triangles or
starting with two overlapping \(K_4\) copies in the one-edge case.

\subsection{6.1 Dominating neighbors of degree five or
six}\label{dominating-neighbors-of-degree-five-or-six}

Let \(v\) have degree 5 and let \(u\) be a neighbor with
\(N[v]\subseteq N[u]\). Write \(W=N(v)-\{u\}\), \(B=G[\{u,v\}\cup W]\),
and \(Z=N(u)-N[v]\). Then \(B\) is \(K_6\) or \(K_6\) minus one rim
edge, and \(|Z|=d(u)-5\).

Take \(X=E(G[W])\cup \{uv\}\cup \{uz:z\in Z\}\). It meets every triangle
at \(u\) or \(v\). If \(B\) is \(K_6-ab\) and \(|Z|\le 1\), \(X\) has at
most 7 edges and \(B\) has the four-packing
\(auv\),\(acd\),\(bud\),\(bcv\) (\(W=\{a,b,c,d\}\)); this is strict. If
\(B=K_6\) and \(Z\) is empty, \(|X|=7\) and again a four-packing is
strict.

If \(B=K_6\) and \(Z=\{p\}\), the edge \(up\) belongs to a triangle, so
\(p\) has a neighbor \(w\) in \(W\) (\(v\) has no neighbor \(p\)).
Choose a four-packing of \(K_6\) avoiding \(uw\) and add \(upw\). Add
\(pw\) to \(X\). Its size is \(9<10\), strict. The required four-packing
exists because \(K_6\) minus a perfect matching decomposes into four
triangles: for nonedge pairs \(ab\),\(cd\),\(st\), use
\(acs\),\(adt\),\(bct\),\(bds\); extend \(uw\) to such a matching.

Thus no dominating neighbor of a degree-five vertex has degree 5 or 6.

\subsection{6.2 Adjacent degree-five
vertices}\label{adjacent-degree-five-vertices}

If two degree-five vertices \(u\),\(v\) are adjacent and have four
common neighbors, they dominate each other and Section 6.1 applies.
Otherwise, their precise links force three common neighbors
\(W=\{w_1,w_2,w_3\}\). Both links are \(K_5\) minus one edge. Let \(p\)
be the unique neighbor of \(u\) outside \(N[v]\) and \(q\) the unique
neighbor of \(v\) outside \(N[u]\). The vertices \(p\),\(q\) are
distinct, each joins all of \(W\), \(vp\) and \(uq\) are absent, and
\(W\) is a triangle. Whether \(pq\) is present is irrelevant.

Use

\[P=\{uw_1w_2, vw_1w_3, upw_3, vqw_2\},\]

\[X=E(W)\cup \{uv,up,vq,pw_3,qw_2\}\]

This is a balanced four/eight reduction deleting \(u\),\(v\). Its
residual \(K\) has equality blocks and \(W\) is independent there. The
five triangles

\[uvw_1, upw_2, vqw_3, pw_1w_3, qw_1w_2\]

are edge-disjoint. Their surviving edges in \(K\) are exactly
\(pw_1\),\(pw_2\),\(qw_1\),\(qw_3\). A residual block meets \(W\) at
most once; it therefore sees at most two of these edges, and any pair
shares that \(W\) vertex. Block avoidance lets us append all five to a
maximum \(K\)-packing, contradiction. Hence the degree-five vertices
form an independent set.

\subsection{6.3 A nondominating degree-six
neighbor}\label{a-nondominating-degree-six-neighbor}

Suppose \(v\) has degree 5 with link \(K_5-ut\), and \(u\) is one of the
missing-edge endpoints. Thus \(N(v)=\{u,t,a,b,c\}\), \(ut\) is absent,
and \(W=\{a,b,c\}\) is a triangle. If \(d(u)=6\), write its two
remaining neighbors \(p\),\(q\) outside \(N[v]\). In \(G-u\), the
vertices \(\{v,t,a,b,c\}\) induce \(K_5\).

If \(pq\) is present, in the link at \(u\) take \(M=\{pq,va,bc\}\) and
\(Q=W\). Outside \(Q\) the only link edge is \(pq\). This balanced
three/six singleton reduction leaves that \(K_5\) minus the two
independent edges \(va\),\(bc\), impossible for residual equality
blocks.

If \(pq\) is absent, triangle codegree at least 2 for \(up\) and \(uq\)
implies that \(p\),\(q\) each have at least two neighbors in \(W\).
Choose distinct a in \(N(p)\cap W\) and \(b\) in \(N(q)\cap W\), with
\(c\) the third vertex. Take \(M=\{pa,qb,vc\}\) and \(Q=W\). The link
outside \(Q\) is independent. This balanced three/six reduction leaves
the \(K_5\) with only \(vc\) deleted, again impossible.

Together with minimum degree 5,6.1 and Section 6.2, this proves:

\textbf{Low-neighbor restriction.} Every neighbor of a degree-five
vertex has degree at least 7.

\section{7. Overlap restriction for nonadjacent degree-five
vertices}\label{overlap-restriction-for-nonadjacent-degree-five-vertices}

\par\medskip\noindent\textbf{Proposition 8 (Overlap restriction; Proposition \(O\)).} Two
nonadjacent degree-five vertices in a minimal combined obstruction have
at most three common neighbors. If they have three, both links are
\(K_5\).

\emph{Proof.}

\textbf{Attribution and extension.} Botler--Fernandes--Gutierrez Lemma
3.3 proves the corresponding implication in irreducible robust graphs:
degrees at most six and a neighborhood union of size at most seven force
both degrees to be five, an intersection of size three, and both links
\(K_5\). For two degree-five vertices, an intersection of at least three
meets that union hypothesis. Its balanced reductions cannot simply
exclude a minimal combined-\(S\) failure. The residual-block arguments
in Sections 7.2--7.4 discharge that extra requirement; this is the
precise distinction between the predecessor and the version needed here.

Write \(C=N(x)\cap N(y)\). The two links are \(K_5\) or \(K_5\) minus
one edge by Section 5. No assumption is made about other degrees.

\subsection{7.1 Five common neighbors}\label{five-common-neighbors}

Their common neighborhood A has five vertices and is \(K_5\) or
\(K_5-ab\). Delete \(x\),\(y\) and \(E(A)\); the cover increment is
\(|E(A)|\).

When \(A=K_5\) on \(\{a,b,c,d,e\}\), the six triangles \(xbc\), \(xde\),
\(ybd\), \(yce\), \(abe\), \(acd\) are disjoint. The cover increment ten
is strictly smaller than twelve.

When \(A=K_5-ab\), omit \(abe\) from that list. The remaining five
triangles exist and the cover increment nine is strictly smaller than
ten.

\subsection{7.2 Four common neighbors}\label{four-common-neighbors}

Write \(W=C\) and let \(p\),\(q\) be the respective unique neighbors
outside the common neighborhood. The graph on \(W\) is \(K_4\) or
\(K_4\) minus one edge. It has two edge-disjoint perfect matchings, so
four triangles, two at each center, form a packing.

Take \(X=E(W)\cup \{xp,yq\}\). It meets all triangles at \(x\),\(y\) and
contains the noncentral edges of this four-packing. If \(W\) has five
edges this is strict, with \(|X|=7\).

If \(W=K_4\), the reduction is balanced four/eight and the residual has
equality blocks with \(W\) independent. The vertex \(p\) has at least
three neighbors in \(W\), because the link at \(x\) misses at most one
edge. Thus there is a page \(p\) \(f\) on some \(W\)-edge \(f\). The
graph on \(\{x,y\}\cup W\) is \(K_6-xy\). Since \(xy\) and \(f\) are
disjoint, it has a four-packing avoiding \(f\). Append the page and
compatible residual block packings, giving \(\nu(G)+1\). Contradiction.

\subsection{\texorpdfstring{7.3 Three common neighbors, a missing edge
inside
\(C\)}{7.3 Three common neighbors, a missing edge inside C}}\label{three-common-neighbors-a-missing-edge-inside-c}

If a common edge \(ac\) is absent, write \(C=\{a,b,c\}\),
\(N(x)-C=\{p,q\}\), \(N(y)-C=\{r,s\}\). Both links are precisely
\(K_5-ac\); all edges required below therefore exist. The pairs
\(p\),\(q\) and \(r,s\) are disjoint, because \(C\) is exactly the
common neighborhood.

Use \(P=\{qab,bcr,xpa,xqc,yar,ycs\}\) and
\(X=E(qab)\cup E(bcr)\cup \{qc,pa,ar,cs,xp,ys\}\). The six triangles are
disjoint and \(|X|=12\). It covers every triangle incident to \(x\) or
\(y\): \(xp\) handles the \(p\)-pages at \(x\), \(qa\),\(qb\),\(qc\)
handle its \(q\)-pages, \(ab\),\(bc\) handle its \(C\)-triangles; the
\(y\) side is analogous. Every noncentral \(P\)-edge lies in \(X\). Thus
the residual has equality blocks.

Now use the seven edge-disjoint triangles \(xpa\), \(qab\), \(pbc\),
\(xqc\), \(yar\), \(ysb\), \(rsc\). Their edges surviving in the
residual are exactly \(pb\),\(pc\),\(sb\),\(rs\). Every pair of these
except \(pb\),\(sb\) requires an absent residual edge to be in a common
clique block: \(pb\)/\(pc\) requires \(bc\), \(sb\)/\(rs\) requires
\(br\), \(pb\)/\(rs\) requires \(br\), \(pc\)/\(sb\) requires \(bc\),
and \(pc\)/\(rs\) requires \(cr\). Those edges were deleted. Hence each
block contains at most two forbidden edges and the only possible pair
shares \(b\). Standard block avoidance appends the seven-packing,
contradiction.

\subsection{\texorpdfstring{7.4 Three common neighbors, a missing edge
outside
\(C\)}{7.4 Three common neighbors, a missing edge outside C}}\label{three-common-neighbors-a-missing-edge-outside-c}

Suppose \(C\) is a triangle and at least one link is \(K_5\) minus an
edge; take that link to be the one at \(y\). Let its missing edge be
\(e\). At least one endpoint of \(e\) lies outside \(C\). Choose an edge
\(f=ac\) of \(C\) disjoint from \(e\) and let \(b\) be the third vertex.

Write \(N(x)-C=\{p,q\}\). At least one of \(p\),\(q\) joins all three
\(C\) vertices; call it \(p\). Relabel \(a\),\(c\) if necessary so that
\(qa\) exists. Then \(pb\) also exists.

The core on \(\{y\}\cup N(y)\) is \(K_6-e\). Because \(e\) and \(f\) are
disjoint, it has a four-packing \(P_0\) avoiding \(f\). Take
\(P=P_0\cup \{xac,xpb\}\) and \(X=E(G[N(y)])\cup \{xp,xq,pb\}\). This is
a balanced six/twelve reduction deleting \(x\),\(y\). The nine edges of
the near-clique link at \(y\) cover all its triangles. Any
\(x\)-triangle either contains \(p\) or \(q\), hit by \(xp\) or \(xq\),
or has its other two vertices in \(C\), whose edges are deleted.
Noncentral packed edges all lie in \(X\).

The residual has equality blocks and \(N(y)\) is independent in it.
Replace the two \(x\)-triangles by three triangles \(pac\), \(xpb\),
\(xqa\). Together with \(P_0\) they form a seven-packing. Its surviving
edges are \(pa\),\(pc\),\(qa\): \(pb\) and \(ac\) were deleted. A
residual block meets \(N(y)\) at most once, so its only possible
two-edge footprint is \(pa\),\(qa\), sharing \(a\). Avoidance gives the
contradiction.

All cases with a near-clique link are excluded, proving Proposition
\(O\).

\section{8. The global peeling reduction and final local
configuration}\label{the-global-peeling-reduction-and-final-local-configuration}

\par\medskip\noindent\textbf{Lemma 9 (Final configuration and exclusion).} A five-degenerate
minimal combined obstruction forces the eight-vertex configuration
described below, with arbitrary outside attachments, and this
configuration cannot occur.

\emph{Proof.}

We now exclude the minimal obstruction in the subgraph-closed class of
all 5-degenerate graphs.

Let \(G\) be a minimal \(S\) failure among 5-degenerate graphs. Let
\(L\) be its set of degree-five vertices. It is nonempty by degeneracy
and Section 4. Section 6 makes \(L\) independent and gives degree at
least 7 for all its neighbors.

The subgraph \(G-L\) is nonempty and 5-degenerate. Choose \(u\) in it
with at most five neighbors in \(G-L\). Set \(A=N(u)-L\), \(d=d(u)\),
\(r=|N(u)\cap L|\). Then \(|A|\le 5\) and \(r\ge d-5\). Since \(u\) is
outside \(L\) and \(\delta(G)\ge 5\), \(d\ge 6\), so \(r\ge 1\). Thus
\(u\) does have a neighbor in \(L\), and the low-neighbor restriction
gives \(d\ge 7\) and \(r\ge 2\).

For each \(v\) in \(N(u)\cap L\), put \(A_v=N(v)\cap A\). If \(u\)
dominates \(v\), \(|A_v|=4\); otherwise \(|A_v|=3\), and \(v\) has one
further neighbor \(p_v\) outside \(N[u]\). These descriptions use the
precise \(K_5\)/\(K_5\)-minus-edge links. For distinct low vertices
\(v\),\(w\), the sets \(A_v\),\(A_w\) intersect in at most one vertex.
Indeed two common A vertices together with \(u\) give at least three
common neighbors. Proposition \(O\) excludes this if either link is a
near-clique. If both links are cliques, both A sets have size four in an
at-most-five-element set, so their intersection has size at least three
and there are at least four common neighbors, again excluded.

There can be at most two subsets of an at-most-five-element set, each of
size at least three, whose pairwise intersections have size at most one.
To see this, two must each have size three, meet in one point and
exhaust five points. Any third three-set would have total intersection
size at least three with their union, whereas meeting each in at most
one allows at most two. Thus \(r=2\), \(d=7\), \(|A|=5\), and both A
sets are triples meeting in exactly one vertex.

Label \(A=\{a,b,c,d,e\}\), \(N(v)=\{u,p,a,b,c\}\),
\(N(w)=\{u,q,a,d,e\}\), \(N(u)=\{v,w,a,b,c,d,e\}\). The low vertices
\(v\),\(w\) have near-clique links missing \(up\),\(uq\) respectively.
The vertices \(p\),\(q\) lie outside \(N[u]\) and are distinct: if they
were equal, \(v\),\(w\) would have three common neighbors
\(u\),\(a\),\(p\) and Proposition \(O\) would exclude them. The sets
\(\{a,b,c\}\) and \(\{a,d,e\}\) are triangles.

\subsection{8.1 The four cross edges are
forced}\label{the-four-cross-edges-are-forced}

Suppose one of the four edges joining \(\{b,c\}\) to \(\{d,e\}\) is
absent. Relabel within each pair so that \(ce\) is absent.

In the link at \(u\), choose \(M=\{vc,we,ab\}\), \(Q=\{a,b,d\}\).
Outside \(Q\) the vertices are \(v\),\(c\),\(w\),\(e\), whose only edges
are \(vc\) and \(we\): \(vw\),\(ve\),\(wc\) are absent by the displayed
neighborhoods and \(ce\) is absent by assumption. Hence this is a
balanced singleton three/six reduction.

The residual contains \(K_5\) on \(\{v,p,a,b,c\}\) with \(vc\) and
\(ab\) deleted. The other deleted edge \(we\) is outside this vertex set
because \(p\) lies outside \(N[u]\). This is impossible for the equality
block classification. Thus all four cross edges exist and \(A\) is
\(K_5\).

\subsection{\texorpdfstring{8.2 The final eight-vertex graph
\(B\)}{8.2 The final eight-vertex graph B}}\label{the-final-eight-vertex-graph-b}

Let \(B\) be the graph induced on \(\{u,v,w,a,b,c,d,e\}\). It has
\(A=K_5\), \(u\) joined to all other seven vertices, \(v\) joined within
\(B\) precisely to \(u\),\(a\),\(b\),\(c\), and \(w\) joined within
\(B\) precisely to \(u\),\(a\),\(d\),\(e\). The only further neighbors
of \(v\),\(w\) in \(G\) are \(p\),\(q\) respectively.

A seven-packing of \(B\) avoiding \(ab\) is \(vua\), \(vbc\), \(wud\),
wae, \(uce\), \(acd\), \(bde\). The two unused \(B\) edges are
\(ub\),\(ab\). By swaps inside \(\{b,c\}\), inside \(\{d,e\}\), and
exchange of the two arms together with \(v\),\(w\), there is a
seven-packing avoiding each \(A\)-edge incident with \(a\).

\textbf{Pair-avoidance claim.} For any two distinct \(A\)-edges
\(f\),\(g\), \(B-\{f,g\}\) has six edge-disjoint triangles.

If one edge is incident to \(a\), take a seven-packing avoiding it and
discard the possible triangle using the other edge. It remains to
consider pairs among the six edges on \(\{b,c,d,e\}\). Up to the same
symmetries their types and witnesses are:

{\def\LTcaptype{none} % do not increment counter
{\small\begin{longtable}[]{@{}
  >{\raggedright\arraybackslash}p{(\linewidth - 2\tabcolsep) * \real{0.5000}}
  >{\raggedright\arraybackslash}p{(\linewidth - 2\tabcolsep) * \real{0.5000}}@{}}
\toprule\noalign{}
\begin{minipage}[b]{\linewidth}\raggedright
Avoided edges / type
\end{minipage} & \begin{minipage}[b]{\linewidth}\raggedright
Six edge-disjoint triangles
\end{minipage} \\
\midrule\noalign{}
\endhead
\bottomrule\noalign{}
\endlastfoot
\(bc\),\(de\): two arm edges & \(vub\), \(vac\), \(wud\), wae, \(uce\),
\(abd\) \\
\(bc\),\(bd\): arm edge and cross edge & \(vub\), \(vac\), \(wue\),
\(wad\), \(ucd\), \(abe\) \\
\(bd\),\(ce\): disjoint cross edges & \(vub\), \(vac\), \(wue\),
\(wad\), \(ucd\), \(abe\) \\
\(bd\),be: adjacent cross edges & \(vua\), \(vbc\), \(wud\), wae,
\(uce\), \(acd\) \\
\end{longtable}}
}

These four cases exhaust pairs: an arm edge and any cross edge meet; two
cross edges either meet or are disjoint. Every displayed triangle uses
only \(B\) edges and the triangles in each row are pairwise
edge-disjoint. Thus this is an ordinary finite certificate with an
explicit symmetry reduction, not an appeal to an unperformed
enumeration.

\subsection{8.3 External-page augmentation closes the
configuration}\label{external-page-augmentation-closes-the-configuration}

Delete \(u\),\(v\),\(w\) and every \(A\)-edge, obtaining \(K\). The
seven-packing of \(B\) appends to every \(K\)-packing. The fourteen
edges \(X=E(A)\cup \{uv,uw,vp,wq\}\) cover every triangle lost in the
deletion. Triangles at \(u\) are inside \(B\); triangles at \(v\) or
\(w\) are inside \(B\) or use \(vp\) or \(wq\); other lost triangles use
an \(A\)-edge. Equality therefore forces \(\nu(K)=\nu(G)-7\),
\(\tau(K)=2\nu(K)\), and \(K\) has \(K_4\)/\(K_5\) triangle components.
Since \(A\) is independent in \(K\), each block meets A at most once.

Let \(F\) be the family of all triangles with an \(A\)-edge as base and
apex outside \(B\). If \(F\) contains two edge-disjoint members on bases
\(f\),\(g\), use a six-packing of \(B\) avoiding \(f\),\(g\) and append
those two pages. Each residual block has at most two old forbidden
edges, sharing its unique A vertex. Maximum residual block packings
avoid them. This constructs \(\nu(K)+8=\nu(G)+1\) triangles,
contradiction.

Thus \(F\) is pairwise edge-intersecting and has a cover of size at most
two by the elementary intersecting-triangle lemma.

The following nine edges cover \(B\):
\(\{ua,ub,ab,cd,ce,de,vc,wd,we\}\). The first six are the within-part
edges of the \(3+3\) partition \(\{u,a,b\}\)/\(\{c,d,e\}\) of the
\(K_6\) core. The only remaining \(v\)-triangles use \(vc\); the only
remaining \(w\)-triangles use \(wd\) or \(we\).

There are four separate cover contributions:

\begin{enumerate}
\def\labelenumi{\arabic{enumi}.}
\tightlist
\item
  A minimum \(K\)-cover handles all triangles surviving in \(K\).
\item
  The displayed nine-edge cover handles all triangles wholly in \(B\).
\item
  An \(F\)-cover of size at most two handles all triangles with an
  \(A\)-edge and apex outside \(B\).
\item
  The two edges \(vp\),\(wq\) handle every remaining triangle incident
  with the deleted low vertices \(v\),\(w\).
\end{enumerate}

To check exhaustiveness, a lost triangle either meets \(\{u,v,w\}\) or
uses a deleted \(A\)-edge. Any triangle at \(u\) is wholly in \(B\)
because \(N(u)=\{v,w\}\cup A\). A triangle at \(v\) either lies in \(B\)
or uses its unique outside-neighbor edge \(vp\); similarly \(w\) and
\(wq\). A triangle using an \(A\)-edge either has its third vertex in
\(B\), or is an external page in \(F\). All graph edges chosen in the
\(F\)-cover are actual edges of \(G\). Contributions can overlap, which
only reduces the size of their union; no disjointness is needed for the
upper bound. Thus a minimum \(K\)-cover plus the other contributions has
size at most \(\tau(K)+9+2+2=\tau(K)+13\), contradicting
\(\tau(G)=\tau(K)+14\).

For the inequality alone, the earlier seven/fourteen reducing tuple
already contradicts irreducibility, as reconstructed in the short
credited proof. The saving from fourteen to thirteen and the argument
for arbitrary outside pages are retained here to exclude equality
obstructions.

Therefore the minimal failure cannot exist.

\section{9. Conclusion, converse and
scope}\label{conclusion-converse-and-scope}

The class of finite 5-degenerate graphs is closed under every vertex and
edge deletion used above. If it contained a failure of \(S\), choose one
with the fewest edges and discard isolates. Sections 1--5 give the
minimal-obstruction premises. Sections 6--8 exclude it. Hence every
graph in this class satisfies the inequality and the asserted necessity
at equality.

Conversely, packing and covering numbers add over triangle components
because their edge supports are disjoint. \(K_4\) has
\((\nu,\tau)=(1,2)\), \(K_5\) has (2,4), and triangle-free edges
contribute zero. Thus the stated component condition implies equality,
including the vacuous triangle-free case.

This proves the main theorem by ordinary finite graph arguments. The
public finite-certificate checkers independently verify displayed local
certificates and the supplied small-link coverage; they do not enumerate
arbitrary external attachments and are not substitutes for the induction
or page arguments above. The formal-verification account appears in
Section 10; the ordinary argument and finite checks have distinct
evidentiary roles. No external human endorsement or established
historical priority is claimed.

The theorem extends the TUZA-04 result by removing its
forbidden-\(K_6\)-minus-edge condition entirely. It does not prove
general Tuza, the global \(K_4\)/\(K_5\) equality conjecture,
unrestricted two-triangle restoration, stable-restoration strictness,
OCT-five, or the fixed-\(c<3\) density application. Those graph classes
need not be 5-degenerate. Each becomes an immediate application only
when the resulting graph is independently 5-degenerate.

\section{10. Formal verification, coverage, and
limitations}\label{formal-verification-coverage-and-limitations}

The exact public target is
\href{https://github.com/pappa211/ai-math-experiments/blob/main/problems/TUZA-06/lean/Tuza06/FiveDegenerate.lean}{\texttt{Tuza06.five\_degenerate\_combined}}:

\begin{Shaded}
\begin{Highlighting}[]
\NormalTok{@Tuza06.five\_degenerate\_combined :}
\NormalTok{  \ensuremath{\forall} \{V : Type u\} [Fintype V] [DecidableEq V] \{G : SimpleGraph V\},}
\NormalTok{    Tuza06.Degenerate G 5 \ensuremath{\rightarrow}}
\NormalTok{      Tuza.tau (Tuza.triangles G) \ensuremath{\le} 2 * Tuza.nu (Tuza.triangles G) \ensuremath{\land}}
\NormalTok{      (Tuza.tau (Tuza.triangles G) = 2 * Tuza.nu (Tuza.triangles G) \ensuremath{\leftrightarrow}}
\NormalTok{        Tuza06.ComponentCliques G)}
\end{Highlighting}
\end{Shaded}

The universe binder above is displayed using the conventional name
\texttt{u}; generated printer names may differ. \texttt{triangles\ G} is
the full \texttt{G.cliqueFinset\ 3}. \texttt{nu} counts edge-disjoint
triangle families. \texttt{tau} uses actual two-vertex edge covers;
\texttt{Tuza.graph\_min\_cover} proves attainment by graph edges.
\texttt{Degenerate} uses the usual induced-subset definition.
\texttt{ComponentCliques} uses reflexive transitive closure of
shared-edge triangle adjacency and literal edge support. The theorem has
no unproved external mathematical inputs. Its only substantive
hypothesis is \texttt{Degenerate\ G\ 5}; the ambient finite-type and
decidable-equality instances are shown explicitly.

The frozen target \texttt{CombinedStatement\ G} is proved by
\texttt{five\_degenerate\_statement}. Definition-stage comments in the
immutable \texttt{Statement.lean} record the earlier target, before
completion; the final theorem and
\href{https://github.com/pappa211/ai-math-experiments/blob/main/problems/TUZA-06/PROOF_STATUS.md}{current proof status}
supersede those historical comments. The four reused \texttt{Tuza}
modules contain unrelated conditional interfaces, but their unproved
mathematical inputs are not premises of this theorem.

{\def\LTcaptype{none} % do not increment counter
{\small\begin{longtable}[]{@{}
  >{\raggedright\arraybackslash}p{(\linewidth - 2\tabcolsep) * \real{0.5000}}
  >{\raggedright\arraybackslash}p{(\linewidth - 2\tabcolsep) * \real{0.5000}}@{}}
\toprule\noalign{}
\begin{minipage}[b]{\linewidth}\raggedright
Paper result or step
\end{minipage} & \begin{minipage}[b]{\linewidth}\raggedright
Actual Lean declaration(s), in namespace \texttt{Tuza06} unless stated
otherwise
\end{minipage} \\
\midrule\noalign{}
\endhead
\bottomrule\noalign{}
\endlastfoot
Theorem 1, both conclusions & \texttt{five\_degenerate\_combined} \\
Separate inequality and equality &
\texttt{five\_degenerate\_inequality};
\texttt{five\_degenerate\_equality} \\
Frozen target and block equivalence &
\texttt{five\_degenerate\_statement};
\texttt{componentCliques\_iff\_decomposition} \\
Equality converse & \texttt{componentCliques\_equality} \\
Lemma 2, clique parameters and avoidance & \texttt{k4\_parameters};
\texttt{k5\_parameters}; \texttt{k5\_avoid\_two};
\texttt{Tuza.k4\_face\_avoiding} \\
Intersecting triangle families & \texttt{Tuza.intersecting\_cover};
\texttt{Tuza.intersecting\_k4\_structure} \\
Lemma 3, minimality and criticality &
\texttt{exists\_minimal\_obstruction};
\texttt{minimal\_tight\_critical} \\
Lemma 4, every edge in two triangles &
\texttt{minimal\_edge\_codegree} \\
Lemma 5, closed link size & \texttt{minimal\_closed\_link\_size} \\
Lemma 6, all degree-five links &
\texttt{minimal\_degree\_five\_configuration} \\
Lemma 7, low adjacency and high neighbors &
\texttt{minimal\_degree\_five\_independent};
\texttt{minimal\_degree\_five\_neighbors\_high} \\
Proposition 8, complete overlap restrictions &
\texttt{minimal\_degree\_five\_common\_le\_three};
\texttt{minimal\_three\_common\_cliques} \\
Lemma 9, second degeneracy and labels &
\texttt{minimal\_two\_arm\_peeling};
\texttt{minimal\_final\_arm\_data} \\
Actual final configuration &
\texttt{final\_configuration\_of\_arm\_data} \\
Arbitrary pages and strict final cover &
\texttt{minimal\_no\_final\_configuration} \\
Auxiliary four-degenerate combined theorem &
\texttt{four\_degenerate\_combined} \\
\end{longtable}}
}

The prescribed-edge minimum-cover observation in Lemma 2 is an
elementary ordinary detail used in Section 4.4. It is not separately
exported as a general \(K_5\) cover declaration. The formal development
instead uses the unified small-link argument in Appendix B; the complete
formal theorem includes that route.

The preserved source freeze covers \textbf{280 mathematical modules},
all in the main entry point's import closure. The source of
\texttt{FiveDegenerate.lean} has SHA-256
\nolinkurl{978ebe0431fde9ecf9598538c6c01e06bcd9ee363bd5ee03a53b76c43143fadf}.
All 280 files are preserved byte-for-byte. Lean \textbf{4.33.1} and
Mathlib revision \nolinkurl{0df444a360eaa60ab8c11dca51a86af692955474},
together with all nine package revisions, are pinned. Complete theorem
types and transitive-axiom audits cover \textbf{618 public theorem
declarations}; the only reported axioms are \texttt{propext},
\texttt{Classical.choice}, and \texttt{Quot.sound}. There is no
\texttt{sorry}, custom mathematical axiom, hidden external theorem
premise, or native proof shortcut in the mathematical closure.

The frozen TUZA-06 verification rebuilt every mathematical module in a
fresh project workspace and replayed all \textbf{280 fresh module
objects} serially through the standard Lean kernel, including private
proof constants. Publication verification separately rebuilds and
replays the exact public source tree; its outcome, commands and evidence
are recorded in the
\href{https://github.com/pappa211/ai-math-experiments/blob/main/problems/TUZA-06/REPRODUCTION.md}{portable reproduction
record}. Reused compiled dependencies were separately pinned and
checked. This is \textbf{not a cold Mathlib rebuild}. The replay adapter
retains the upstream LeanChecker replay routine; its unsafe
IO/runtime-memory wrapper is outside the mathematical import closure.
Replay is additional checking by the same standard Lean kernel,
\textbf{not an independently implemented kernel}. Upstream source
provenance and license accompany the tool.

The independent Python checks validate supplied local finite
certificates, five-link coverage and overlap witnesses. They do not
enumerate arbitrary outside attachments and are not premises of the
formal theorem. The symbolic page and compatibility arguments in Section
8 discharge that unbounded obligation. Neither kernel checking nor
internal review certifies novelty, human specialist endorsement, or
broader conjectures.

The public \href{https://github.com/pappa211/ai-math-experiments/blob/main/problems/TUZA-06/README.md}{research archive}
supplies exact Lean sources, full public proof records, relevant finite
certificates/checkers, attribution and provenance, selected reviews,
historical TUZA-04/05 proofs and limitations, source-to-publication
mapping, and reproduction commands. These records distinguish earlier
partial status from the completed theorem. Operational instructions and
private review archives are excluded.

\section{Appendix A. A short credited proof of the
inequality}\label{appendix-a.-a-short-credited-proof-of-the-inequality}

\subsection{\texorpdfstring{\(A\).1 Published inputs and exact
scope}{A.1 Published inputs and exact scope}}\label{a.1-published-inputs-and-exact-scope}

Puleo, {[}1{]}, Definition 2.1 and Lemma 2.2, supplies the reducible-set
framework. Lemma 2.6 gives robustness for a minimal inequality
counterexample. In a robust graph with no reducible set, Lemma
2.7(a),(\(b\)), equivalently Propositions 5.1 and 5.2, implies that
degree-five links are \(K_5\) or \(K_5\) minus one edge, and vertices of
degree at most six form an independent set. Here robustness means that
every component of every vertex link has at least five vertices.
Complements in Proposition 5.1 are essential; plain PDF text extraction
can lose their overbars.

Botler--Fernandes--Gutierrez, {[}2{]}, Lemma 3.3, applies to irreducible
robust graphs. For two degree-five vertices with at least three common
neighbors it forces exactly three common neighbors and both links
\(K_5\). This is the published antecedent used below. Their theorem
concerns treewidth; it is this local lemma that is used here.

\subsection{\texorpdfstring{\(A\).2 Reduction
accounting}{A.2 Reduction accounting}}\label{a.2-reduction-accounting}

A nonempty vertex set \(U\) is reducible if there are \(r\)
edge-disjoint triangles \(P\) and a graph-edge set \(X\), with
\(|X| \le  2r\), such that \(X\) hits every triangle incident with \(U\)
and contains all \(P\)-edges with both ends outside \(U\). Put
\(K\)=(\(G-X\))-\(U\). Any packing of \(K\) combines with \(P\), and any
\(K\)-cover combines with \(X\) to cover \(G\). Therefore

\[\nu(G) \ge  \nu(K)+r,\]
\[\tau(G) \le  \tau(K)+|X| \le  2\nu(K)+2r \le  2\nu(G)\]

whenever \(\tau(K) \le  2\nu(K)\). The balance \(|X|=2r\) is sufficient;
no strict saving is required for the inequality.

Assume for a contradiction that a 5-degenerate counterexample exists.
Choose one with the fewest edges, discard isolated vertices, and, if
necessary, minimize the number of vertices among equal edge counts.
Every proper subgraph satisfies the inequality: it is still
5-degenerate. Hence the published robustness lemma applies, and no
reducible vertex set exists. Denote this graph by \(G\). Robustness
gives minimum degree at least five. The published low-degree facts above
give:

\begin{itemize}
\tightlist
\item
  every degree-five vertex has link \(K_5\) or \(K_5\) minus one edge;
\item
  degree-five vertices are independent and all their neighbors have
  degree at least seven;
\item
  two degree-five vertices have at most three common neighbors, and if
  they have three then both links are cliques.
\end{itemize}

These are uses of inequality irreducibility. They are not asserted for a
combined equality obstruction without further proof.

\subsection{\texorpdfstring{\(A\).3 Degeneracy
peeling}{A.3 Degeneracy peeling}}\label{a.3-degeneracy-peeling}

Let \(L\) be the set of degree-five vertices. It is nonempty by
5-degeneracy and minimum degree five. It is independent. \(G-L\) is
nonempty because a vertex in \(L\) has neighbors outside \(L\). Choose
\(u\) in \(G-L\) with degree at most five in \(G-L\), and set

\[A=N(u)-L,  r=|N(u)\cap L|,  d=d_G(u)\]

Then \(|A|\le 5\), \(r=d-|A|\ge d-5\), and \(d\ge 6\) because \(u\) is
outside \(L\). Thus \(r\ge 1\). A neighbor in \(L\) forces \(d\ge 7\),
so \(r\ge 2\). The second use of degeneracy bounds degree in \(G-L\),
not degree in \(G\).

For each \(v\) in \(N(u)\cap L\), let \(A_v=N(v)\cap A\). Because \(L\)
is independent, every neighbor of \(v\) other than \(u\) lies outside
\(L\). If the link at \(v\) is a clique, \(u\) dominates \(v\) and
\(|A_v|=4\). If that link misses an edge not incident with \(u\), the
same conclusion holds. If its missing edge is \(up_v\), then \(|A_v|=3\)
and \(p_v\) is the unique neighbor of \(v\) outside \(N[u]\). These
exhaust the cases. In all cases \(|A_v|\ge 3\).

For distinct low neighbors \(v\),\(w\) of \(u\), suppose
\(|A_v\cap A_w|\ge 2\). The common neighbors of \(v\),\(w\) include
\(u\) and those two vertices, so the published overlap lemma makes both
links cliques and gives exactly three common neighbors. But then \(A_v\)
and \(A_w\) are both four-subsets of the at-most-five-element set \(A\),
so they meet in at least three vertices. Together with \(u\) this gives
at least four common neighbors, a contradiction. Hence the sets \(A_v\)
have pairwise intersections at most one.

There are at most two subsets of a set of size at most five, each of
size at least three and pairwise intersecting in at most one element.
Indeed, any two must be three-sets meeting in one point and covering all
five points; a third three-set cannot meet each of the first two in at
most one point. Since \(r\ge 2\), it follows that \(r=2\), \(d=7\),
\(|A|=5\), and the two sets \(A_v\),\(A_w\) are triples meeting in one
point.

Label

\[A=\{a,b,c,d,e\},\] \[N(v)=\{u,p,a,b,c\},\] \[N(w)=\{u,q,a,d,e\},\]
\[N(u)=\{v,w,a,b,c,d,e\}\]

The overloaded letter \(d\) in this display is now a vertex label; the
degree of \(u\) is seven. Both low links are near-cliques, missing
\(up\) and \(uq\) respectively. The vertices \(p\),\(q\) lie outside
\(N[u]\). They are distinct, because otherwise \(v\),\(w\) would have
common neighbors \(u\),\(a\),\(p\) while their links were not cliques.
The sets \(\{a,b,c\}\) and \(\{a,d,e\}\) are triangles, and \(p\) joins
\(a\),\(b\),\(c\) while \(q\) joins \(a\),\(d\),\(e\).

\subsection{\texorpdfstring{\(A\).4 The final
reduction}{A.4 The final reduction}}\label{a.4-the-final-reduction}

If a cross edge between \(\{b,c\}\) and \(\{d,e\}\) is missing, swap
labels within the two pairs so that \(ce\) is absent. In the link at
\(u\) use the matching

\[M=\{vc,we,ab\}\]

and \(Q=\{a,b,d\}\). The link after removing \(Q\) has vertices
\(v\),\(c\),\(w\),\(e\) and only edges \(vc\),\(we\): the other possible
edges are absent from the displayed neighborhoods or by \(ce\) being
missing. Therefore the three triangles

\[uvc, uwe, uab\]

and the six edges

\[X=\{vc,we,ab,ua,ub,ud\}\]

make \(\{u\}\) reducible. Every \(u\)-triangle either has a neighbor in
\(Q\), or uses a matching edge outside \(Q\). The noncentral edges of
all three packed triangles lie in \(X\). This contradicts
irreducibility. Thus all four cross edges exist and \(A\) is \(K_5\).

Let \(B\) be the local graph on \(\{u,v,w,a,b,c,d,e\}\). Its edges are
the ten \(A\)-edges, all seven edges from \(u\) to the other local
vertices, \(va\),\(vb\),\(vc\), and \(wa\),\(wd\),\(we\). The following
seven triangles exist and are pairwise edge-disjoint:

\[vua, vbc, wud, wae, uce, acd, bde\]

Take \(U=\{u,v,w\}\) and

\[X=E(A)\cup \{uv,uw,vp,wq\}\]

There are fourteen edges in \(X\). Every triangle at \(u\) is wholly in
\(B\) and either uses \(uv\), \(uw\), or an \(A\)-edge. Every triangle
at \(v\) or \(w\) either uses the corresponding outside-neighbor edge
\(vp\) or \(wq\), or is local; its local case uses \(uv\), \(uw\), or an
\(A\)-edge. Thus \(X\) meets every triangle incident with \(U\). Every
edge of the displayed seven-packing outside \(U\) is an \(A\)-edge and
belongs to \(X\). Hence this is a balanced seven/fourteen reduction,
again contradicting irreducibility.

No minimal inequality counterexample exists, proving
\(\tau(G)\le 2\nu(G)\) for every finite simple 5-degenerate graph \(G\).
Empty and triangle-free graphs have both parameters zero and satisfy the
assertion directly.

\section{Appendix B. The small-link proof used by the formal
development}\label{appendix-b.-the-small-link-proof-used-by-the-formal-development}

The formal development unifies the diamond and four-cycle arguments of
Sections 4.3--4.4 as follows. It changes the organization, not the
theorem or attribution boundary.

Let \(C=\{a,b,c,d\}\) be closed in the link at \(v\), with \(ab\) absent
and \(ac\),\(ad\),\(bc\),\(bd\) present; \(cd\) is unrestricted. Delete
only the four \(v\)-spokes into \(C\) and rim edges \(ac\),\(bd\). The
two triangles \(vac\),\(vbd\) append to every residual packing, and
\(ac\),\(bd\),\(vc\),\(vd\) cover every lost triangle. Balanced two/four
accounting gives residual equality blocks and the exact increment
\(\nu(G)=\nu(H)+2\). Closure restricts only link edges; arbitrary
external pages remain allowed at this point.

For an external \(ac\)-page with apex \(x\), combine that page with
\(vad\),\(vbc\). Their surviving footprint is
\(ad\),\(bc\),\(xa\),\(xc\). Since \(ac\) is absent in \(H\), the two
terminals \(a\),\(c\) are independent; each residual clique meets at
most one terminal and sees at most two forbidden incident edges. Maximum
\(K_4\)/\(K_5\) block packings avoid this footprint, giving
\(\nu(H)+3\), a contradiction. The \(bd\) case is the same explicit
relabeling with terminals \(b\),\(d\). Thus every lost triangle using
\(ac\) or \(bd\) is local. The remaining lost hub triangles are covered
by \(vc\),\(vd\). Every rim triangle must contain \(c\),\(d\) because
\(ab\) is absent. If \(cd\) exists, adding \(cd\) yields a three-edge
lost-triangle cover; if \(cd\) is absent, there is no rim triangle and
\(vc\),\(vd\) suffice. Either contradicts the exact four-edge increment.
This excludes both local link types.

Other formal changes are organizational: finite label coverage is proved
explicitly; all reductions are instantiated from actual neighborhoods;
outside pages and residual block compatibility are handled symbolically.
Their exact declarations, dependencies and completion status belong to
the formal status record, rather than to an inference from this
manuscript.

\section{References}\label{references}

{[}1{]} Gregory J. Puleo. \emph{Tuza's conjecture for graphs with
maximum average degree less than 7}. European Journal of Combinatorics
\textbf{49} (2015), 134--152.
\href{https://doi.org/10.1016/j.ejc.2015.03.006}{DOI:10.1016/j.ejc.2015.03.006}.
Numbering checked against the
\href{https://arxiv.org/pdf/1308.2211v3}{full arXiv version 3}, 13 April
2015.

{[}2{]} Fábio Botler, Cristina G. Fernandes, and Juan Gutiérrez.
\emph{On Tuza's conjecture for triangulations and graphs with small
treewidth}. Discrete Mathematics \textbf{344} (2021), article 112281.
\href{https://doi.org/10.1016/j.disc.2020.112281}{DOI:10.1016/j.disc.2020.112281}.
Local-lemma numbering checked against the
\href{https://arxiv.org/pdf/2002.07925v2}{full arXiv version 2}, 16 July
2020.

{[}3{]} Qing Cui, Penny Haxell, and Will Ma. \emph{Packing and covering
triangles in planar graphs}. Graphs and Combinatorics \textbf{25}
(2009), 817--824; published online 10 February 2010.
\href{https://doi.org/10.1007/s00373-010-0881-5}{DOI:10.1007/s00373-010-0881-5}.
\href{https://www.columbia.edu/~wm2428/papers/packing_covering_triangles.pdf}{Author-hosted
full text}.

{[}4{]} AI Math Experiments. \emph{TUZA-04--06 public mathematical and
formal-verification record}. 30 October 2026.
\href{https://github.com/pappa211/ai-math-experiments/blob/main/problems/TUZA-06/README.md}{Research archive and source
inventory}, \href{https://github.com/pappa211/ai-math-experiments/blob/main/problems/TUZA-06/FORMAL_STATEMENT.md}{formal
target}, and \href{https://github.com/pappa211/ai-math-experiments/blob/main/problems/TUZA-06/REPRODUCTION.md}{portable
reproduction}. Internal AI research and review records; no external
peer-review status is implied.

{[}5{]} AI Math Experiments. \emph{Triangle Restoration and
Packing--Covering Stability in Tuza's Conjecture}. 29 October 2026.
\href{https://github.com/pappa211/ai-math-experiments/blob/main/papers/tuza-triangle-restoration/main.md}{Separate canonical paper}.

The bibliography is also supplied as \href{https://github.com/pappa211/ai-math-experiments/blob/main/papers/tuza-five-degenerate/references.bib}{BibTeX}.
Third-party full texts were consulted for attribution and are linked
rather than redistributed. Mathematical facts are not claimed as
property of the project.

\end{document}
